\documentclass[12pt]{article}
\usepackage[margin=1in]{geometry}
\usepackage{fancyhdr}
% Using default LaTeX font (Computer Modern) - times package fonts not available
\usepackage{parskip} % Handles paragraph spacing

% --- HEADER CONFIGURATION ---
\pagestyle{fancy}
\fancyhf{} % Clear default header/footer
\setlength{\headheight}{26pt} % Increase height to fit 2 lines of text

% Left Header: 2 lines
\lhead{Statement of Purpose\\Carnegie Mellon University} 

% Right Header: Name and Page Number
\rhead{Harsh Gujarathi \quad \thepage}
\setlength{\headsep}{12pt} % Adjust this value to make it smaller or larger
\begin{document}

% --- BODY CONTENT ---

% Introduction (No bold header, standard practice)
While the world focuses on building larger Large Language Models (LLMs), a critical gap in training and serving them is becoming increasingly glaring and undeniable. The next big breakthrough in AI will not just be a more capable model; it will be engineering innovation that makes these models cost-effective and scalable for real-world industries. Driven by this motivation, I seek to pursue a \textbf{Master degree in Artificial Intelligence and Innovation} at Carnegie Mellon University’s Language Technologies Institute and work at the \textbf{intersection of Distributed Systems and Machine Learning}. My undergraduate studies in Computer Science at BITS Pilani, where I was awarded a Merit Scholar for ranking in the \textbf{top 2 percent} out of 850 students, followed by my role as Software Engineer at Microsoft, have prepared me well for this challenge.

% Undergraduate Research (IoT)
\vspace{0.5em}
My first encounter with resource-constrained systems began during my \textbf{undergraduate years at the IoT Research Lab} at BITS under Prof. Neena Goveas. As a student developer, my primary responsibility was to build \textbf{Divyang App}, an Android application for a social cause at the Rotary Club of Mumbai that automated their process of accessing and recommending prosthetics for amputees using the device cameras and sensors. We identified the user needs and successfully deployed an end-to-end ML powered solution with a \textbf{MobileNet SSD finetuned in tensorflow lite}. This demonstrated that delivering tangible value was a greater priority than maximizing a raw accuracy metric, thereby shifting focus to computational efficiency to ensure the solution was accessible and functional for all users.

% Research (Solar)
\vspace{0.5em}
In another research-oriented project with Prof. Neena, we worked on creating a reproducible digital twin for solar panel installations. To simulate the panels we benchmarked multiple approaches, from classic gradient boosting techniques like CatBoost, LightGBM and XGBoost to complex Physics-Informed Neural Networks (PINNs). The classic ensembles already delivered a robust \textbf{89.5\% accuracy}. While my PINN implementation delivered a \textbf{10\% increase} in accuracy over the classic models, it incurred a \textbf{5x penalty} in model size.  For a Digital Twin designed for real-time simulation, this penalty was too prohibitive. This meaningful trade-off between fidelity and efficiency was an eye-opener regarding the \textbf{inherent constraints of physical systems}. This realization cemented my focus on ML Systems, driving me to build practical, scalable solutions that operationalize advancing research.

% Applied AI (DevRev) - Replaces the Zulip/Privacy Paragraph
\vspace{0.5em}
Seeking to translate this focus on efficiency from research prototypes to production-grade systems, I interned at DevRev, an AI-native startup, where I reduced user's notification load by \textbf{6x} and directly boosted notification consumption rates by solving the critical issue of 'notification fatigue.' To achieve this, I architected a hierarchical clustering engine that grouped high-volume streams (batches of $\sim$50) based on semantic parameters like topic and sender. I then integrated a \textbf{GPT-4} and \textbf{LangChain} pipeline to synthesize these clusters into concise, actionable summaries, effectively transforming raw logs into intelligent insights. Furthermore, I ensured system resilience during peak traffic by implementing \textbf{smart bucket-based throttling} to dynamically prioritize critical alerts, maintaining responsiveness while handling global scale.

% % Research (Privacy/Zulip)
% \vspace{0.5em}
% To further explore the breadth of computer science applications, I pursued a project with Prof. Kunal, where we engineered a privacy-preserving middleware layer atop the open-source messaging Zulip architecture to address metadata leakage. We deployed a distributed testbed comprising a central server node and distinct client nodes, developing a custom wrapper to intercept and reroute traffic through \textbf{mixnet-style message batching} and \textbf{encrypted sender identifiers}. This architecture successfully anonymized sender identity and mitigated \textbf{trace-through traffic analysis}, within our controlled experimental environment. However, moving to an open-world deployment exposed a critical vulnerability; without strict configuration enforcement, user-driven modifications render the encrypted payloads \textbf{uninterpretable by the server}, resulting in message failures. This infrastructure reliability issue highlights a major gap between deterministic environments and unpredictable real-world deployments. This experience pushes me to architect resilient distributed systems that ensure complex AI workloads remain robust and scalable in production.

% Work Experience
\vspace{0.5em}
While various projects strengthened my desire to pursue academic excellence, I also wanted to experience software on a large scale. So, after graduation, I joined Microsoft as a Software Engineer, where I engineered a critical fallback extension that drove assessment coverage to near-perfect \textbf{99.9\%} across dataset of over \textbf{100K }enterprise\textbf{ workloads}. I also designed a comprehensive test layer that covers end-to-end scenarios for a service spanning \textbf{36 geographies}, ensuring consistent reliability across diverse global regions. However, integrating our AppCAT analysis tool revealed a staggering reality: modernizing a single legacy application required an average of 10 engineering hours. Accounting to over \textbf{1 million human hours} for modernization, the data proved that manual refactoring is unsustainable at an enterprise scale. Driven by this inefficiency, I envisioned the automatic modernization of web applications, an insight that has sparked my curiosity on the potential use of LLMs to transform legacy systems, which I aim to explore in my future work.

% Leadership
\vspace{0.5em}
Apart from research and engineering, communication and leadership have also been equally important in shaping my overall identity. Throughout my four years at BITS, I served as a \textbf{Teaching Assistant }for\textbf{ seven courses}, ranging from \textbf{core} Operating Systems to \textbf{graduate-level} Software for Embedded Systems. This experience honed my ability to distill complex technical concepts and deliver them at scale, making intricate subjects accessible and engaging. My leadership experience is grounded in resource management; as the \textbf{Undergraduate Lab Coordinator} for the IoT Research Lab, I was responsible for ensuring equitable access to hardware and managing the logistics of multiple student projects. Furthermore, as the \textbf{Academic Head} of the Computer Science Association, I took the initiative to establish new communication channels between students and faculty, organizing engagement sessions to align academic goals with industry and research trends.

% Why CMU
\vspace{0.5em}
Large-scale systems to handle machine learning workloads are rapidly shaping the technical evolution trajectory; however, the rising costs of training and deploying modern models reduce their accessibility. I strive to contribute to this ever-evolving space of distributed systems and machine learning that makes workloads optimized, scalable, and sustainable. The MSAII program at CMU is well suited for my goal of becoming a technical leader who can translate AI capabilities into deployable, scalable products. To bridge the gap between my systems engineering background and advanced model development, I am eager to engage with the rigorous curriculum at the Language Technologies Institute. I specifically intend to take \textbf{Introduction to Deep Learning} to gain the algorithmic intuition required to optimize model architectures, and \textbf{Advanced NLP}, which will provide the theoretical backing for my work on semantic code analysis. I am particularly inspired by the work of \textbf{Professor Graham Neubig} on code generation and evaluation, which directly aligns with my vision of automating legacy software modernization. Additionally, the research of \textbf{Professor Tianqi Chen} on machine learning compilation and efficient deployment resonates deeply with my drive to solve the "last mile" problem of AI cost and latency. I plan to synthesize these technical insights into the \textbf{Artificial Intelligence Innovation Capstone}, transforming my prototype for automated technical debt reduction into a robust, commercially viable platform.

% Conclusion
\vspace{0.5em}
\noindent Thus, my journey has built a solid foundation for me to pursue my dreams. With an enriching exposure to a vast set of challenges and experiences, I am confident that I have the necessary skill set and motivation to pursue the graduate program at CMU. With its strong focus on language technologies and innovation, CMU offers the ideal environment to build the technical depth I require. Lastly, I hope to leverage the expertise I develop during my graduate studies to build a \textbf{venture that delivers optimized computing solutions, helping democratize access to advanced machine learning by reducing the massive compute costs that currently constrain the field}.

\end{document}